\documentclass{article}
\usepackage{fontspec}
\setmainfont{Times New Roman}
\usepackage{amsmath}
\usepackage{amssymb}
\usepackage{amsfonts}

\begin{document}

\title{Лекционни записки: Основи на вероятностите}
\author{}
\date{}
\maketitle

\section{Въведение. Дефиниция за вероятност}

Днес ще продължим с това да дефинираме строго що е то вероятност. Започваме със строгата дефиниция. Нека имаме едно множество от елементарни събития $\Omega$. Това са всичките възможни изходи на някакъв експеримент. Към него разглеждаме и една $\sigma$-алгебра $\mathcal{A}$ (колекция от подмножества на $\Omega$). Най-често, ако говорим за дискретни вероятности, това може да са всички подмножества на $\Omega$.

Тогава функцията $\mathbb{P}: \mathcal{A} \to [0, 1]$ се нарича \textbf{вероятностна функция} (или просто \textbf{вероятност}), ако са изпълнени следните три условия (аксиоми):

\begin{enumerate}
    \item \textbf{Вероятност на достоверното събитие:} 
    \[ \mathbb{P}(\Omega) = 1 \]
    Тъй като $\Omega$ винаги принадлежи на $\sigma$-алгебрата, това свойство по съвсем естествен начин ни казва, че вероятността да се падне някое от всички възможни елементарни събития (т.е. експериментът да завърши с някакъв резултат) е точно $1$.

    \item \textbf{Вероятност на противоположното събитие:} Ако събитието $A$ принадлежи на $\sigma$-алгебрата (т.е. $A \in \mathcal{A}$), то допълнението му $A^c$ също принадлежи на $\mathcal{A}$, и вероятността на допълнението е:
    \[ \mathbb{P}(A^c) = 1 - \mathbb{P}(A) \]

    \item \textbf{Свойство за изброима адитивност:} Ако имаме редица от събития $A_1, A_2, \dots$ от $\sigma$-алгебрата, които са две по две несечещи се (т.е. $A_i \cap A_j = \emptyset$ за $i \neq j$), то вероятността на тяхното обединение е сумата от вероятностите им:
    \[ \mathbb{P}\left(\bigcup_{i=1}^{\infty} A_i\right) = \sum_{i=1}^{\infty} \mathbb{P}(A_i) \]
\end{enumerate}

В този смисъл вероятността отговаря на класическото разбиране, което имаме за мярка. Ако искаме да пресметнем общото тегло на всички нас, можем да сумираме теглата на всеки от нас индивидуално, защото телата ни не се пресичат. Аналогично, ако искаме да пресметнем лицето на част от пода, можем да сумираме плочките, които попадат в тази част, стига те да не се припокриват. Грубо казано, вероятността на изброимо много несечещи се събития е сумата от индивидуалните вероятности. Това е вероятността чисто като мярка – тя удовлетворява тези свойства, които я правят вероятност.

Преди да преминем към следствията, нека въведем една операция над множества. Ако $A$ и $B$ са две събития от $\mathcal{A}$, разликата $A \setminus B$ разбираме като:
\[ A \setminus B = A \cap B^c \]
Това са всички елементарни събития, които попадат в $A$, но не попадат в $B$. Геометрично, ако начертаем множествата $A$ и $B$, $A \setminus B$ е сечението на $A$ с допълнението на $B$.

\section{Следствия от дефиницията за вероятност}

Нека имаме $\sigma$-алгебра $\mathcal{A}$ и вероятност $\mathbb{P}$. В сила са следните твърдения:

\subsection{а) Вероятност на празното множество}
\[ \mathbb{P}(\emptyset) = 0 \]
\textbf{Доказателство:} Празното множество е допълнението на пълното множество $\Omega$ (т.е. $\emptyset = \Omega^c$). Използвайки второто свойство:
\[ \mathbb{P}(\emptyset) = \mathbb{P}(\Omega^c) = 1 - \mathbb{P}(\Omega) = 1 - 1 = 0 \]
Това е интуитивно ясно – вероятността да не се случи нищо е нула, тъй като все нещо трябва да се случи при експеримента.

\subsection{б) Разбиване на събитие}
За всеки $A, B \in \mathcal{A}$:
\[ \mathbb{P}(B) = \mathbb{P}(A \cap B) + \mathbb{P}(A^c \cap B) \]
\textbf{Доказателство:} Всяко множество $B$ може да се представи като обединение на две несечещи се множества: частта на $B$, която попада в $A$, и частта на $B$, която е извън $A$. Тоест:
\[ B = (A \cap B) \cup (A^c \cap B) \]
Тъй като $(A \cap B)$ и $(A^c \cap B)$ не се пресичат, можем да приложим аксиомата за събираемост, получавайки точно търсената формула.

\subsection{в) Монотонност на вероятността}
Ако $A \subseteq B$, то $\mathbb{P}(A) \le \mathbb{P}(B)$.
\textbf{Доказателство:} Щом $A \subseteq B$, можем да представим $B$ като:
\[ B = A \cup (A^c \cap B) \]
Като приложим свойството от точка б) и вземем предвид, че сечението $A \cap B = A$ (защото $A$ е подмножество на $B$), получаваме:
\[ \mathbb{P}(B) = \mathbb{P}(A) + \mathbb{P}(A^c \cap B) \]
Тъй като $\mathbb{P}(A^c \cap B) \ge 0$, следва, че:
\[ \mathbb{P}(B) \ge \mathbb{P}(A) \]

Оттук лесно можем да изведем и формулата за вероятност на обединение на две събития. Тъй като $A \cup B = A \cup (A^c \cap B)$ (където $A$ и $A^c \cap B$ не се пресичат), то:
\[ \mathbb{P}(A \cup B) = \mathbb{P}(A) + \mathbb{P}(A^c \cap B) \]
Замествайки $\mathbb{P}(A^c \cap B) = \mathbb{P}(B) - \mathbb{P}(A \cap B)$ (от точка б), получаваме класическото:
\[ \mathbb{P}(A \cup B) = \mathbb{P}(A) + \mathbb{P}(B) - \mathbb{P}(A \cap B) \]
Ако събитията са взаимно изключващи се ($A \cap B = \emptyset$), членът $\mathbb{P}(A \cap B) = 0$ и вероятността на обединението е просто сумата им.

\subsection{г) Непрекъснатост в нулата}
Ако имаме намаляваща редица от събития $A_1 \supseteq A_2 \supseteq \dots \supseteq A_n \supseteq \dots$, такива че тяхното сечение е празното множество ($\bigcap_{j=1}^\infty A_j = \emptyset$), то границата на техните вероятности е нула:
\[ \lim_{j \to \infty} \mathbb{P}(A_j) = 0 \]
\textbf{Доказателство:} Идеята е да изразим множеството $A_j$ като изброимо обединение на несечещи се "пръстени" (венци). Можем да представим $A_j$ като:
\[ A_j = (A_j \setminus A_{j+1}) \cup (A_{j+1} \setminus A_{j+2}) \cup \dots = \bigcup_{n=j}^{\infty} (A_n \setminus A_{n+1}) \]
Тези разлики $A_n \setminus A_{n+1}$ са взаимно несечещи се по построение. Поради адитивността на вероятността имаме:
\[ \mathbb{P}(A_j) = \sum_{n=j}^{\infty} \mathbb{P}(A_n \setminus A_{n+1}) \]
За $j=1$ получаваме:
\[ \mathbb{P}(A_1) = \sum_{n=1}^{\infty} \mathbb{P}(A_n \setminus A_{n+1}) \le 1 \]
Тъй като редът с общ член $\mathbb{P}(A_n \setminus A_{n+1})$ е сходящ (сумата му е ограничена от $1$), то остатъчната сума на този сходящ ред трябва да клони към нула, когато $j$ клони към безкрайност. Именно тази остатъчна сума е равна на $\mathbb{P}(A_j)$. Следователно $\lim_{j \to \infty} \mathbb{P}(A_j) = 0$. Този резултат е много важен и се нарича "непрекъснатост в нулата".

\subsection{д) Полуадитивност (субадитивност)}
Ако $A_1, A_2, \dots$ са произволни събития (не задължително несечещи се), то вероятността на тяхното обединение е по-малка или равна на сумата от техните вероятности:
\[ \mathbb{P}\left(\bigcup_{j=1}^{\infty} A_j\right) \le \sum_{j=1}^{\infty} \mathbb{P}(A_j) \]
Идеята на доказателството е да конструираме нова редица от непресичащи се събития $C_1, C_2, \dots$, които покриват същото обединение. Полагаме $C_1 = A_1$, $C_2 = A_2 \setminus A_1$, $C_3 = A_3 \setminus (A_1 \cup A_2)$ и общо:
\[ C_n = A_n \setminus \bigcup_{i=1}^{n-1} A_i \]
Тези $C_j$ са непресичащи се и $\bigcup C_j = \bigcup A_j$. Тъй като $C_j \subseteq A_j$, от монотонността (точка в) имаме $\mathbb{P}(C_j) \le \mathbb{P}(A_j)$. Тогава:
\[ \mathbb{P}\left(\bigcup A_j\right) = \mathbb{P}\left(\bigcup C_j\right) = \sum \mathbb{P}(C_j) \le \sum \mathbb{P}(A_j) \]

\section{Дискретна вероятност}

Нека сега разгледаме как конкретно се задава вероятност, като започнем с дискретната.
Най-простият прототип е експеримент с два възможни изхода, например 0 и 1 (ези/тура, успех/неуспех).
Тук $\Omega = \{0, 1\}$. Дефинираме вероятността да се падне $1$ като $\mathbb{P}(\{1\}) = p$ (за някакво число $0 \le p \le 1$), а вероятността да се падне $0$ е $\mathbb{P}(\{0\}) = 1 - p$. Често бележим $q = 1 - p$.

\subsection{Краен брой елементарни събития}
Нека $\Omega = \{1, 2, \dots, N\}$. Това кодира всеки експеримент с $N$ възможни изхода.
Нека са ни дадени $N$ на брой неотрицателни числа $p_1, p_2, \dots, p_N$, такива че:
\[ \sum_{i=1}^N p_i = 1 \]
За всяко събитие $A \subseteq \Omega$ дефинираме вероятността като:
\[ \mathbb{P}(A) := \sum_{i \in A} p_i \]
В частност, вероятността на единичен изход е $\mathbb{P}(\{i\}) = p_i$.
Тази функция реално е вероятност. Сумата $\sum_{i \in A} p_i$ винаги е между 0 и 1. Също така, аксиомата за противоположното събитие е изпълнена:
\[ \mathbb{P}(A^c) = \sum_{i \in A^c} p_i = 1 - \sum_{i \in A} p_i = 1 - \mathbb{P}(A) \]

\subsection{Равномерна вероятност}
Равномерната вероятност е частен случай, при който всички $p_i$ са равни:
\[ p_i = \frac{1}{N}, \quad 1 \le i \le N \]
Този модел се използва за честни игри, например при хвърляне на зар или рулетка, или в байесовата статистика като неинформативно априорно разпределение, когато нямаме допълнителна информация за това кой изход е по-вероятен.

\textbf{Пример (Тото 6/49):} 
Броят на възможните шестици е $N = 13983816$. Ако играта е честна, вероятността да се падне коя да е конкретна шестица (например $\{1, 2, 3, 4, 5, 6\}$) е точно същата като всяка друга:
\[ \mathbb{P}(\{\text{конкретна шестица}\}) = \frac{1}{13983816} \]
Много хора играят с числа от 1 до 31 (поради дати на раждане и т.н.). Вероятността тези числа да се паднат не е по-голяма, но ако се паднат "по-големи" числа, има по-малко участници, с които ще се раздели печалбата. Самата вероятност за печалба обаче е напълно фиксирана и униформна.

\subsection{Изброимо множество от елементарни събития}
Често е удобно пространството $\Omega$ да бъде безкрайно, но изброимо, например $\Omega = \{1, 2, 3, \dots\}$ или $\{0, 1, 2, \dots\}$. Това се прилага, когато моделираме брой събития, време до повреда (дискретизирано) и др. Не е нужно да слагаме изкуствена горна граница (например 10 000 дни за процесор), тъй като математически е по-удобно да допуснем потенциална безкрайност.

Тук имаме редица от вероятности $p_i \ge 0$, за които:
\[ \sum_{i=1}^{\infty} p_i = 1 \]
Важен факт: \textbf{Не можем да дефинираме равномерна вероятност върху изброимо безкрайно множество.} Ако всички $p_i$ са равни на някакво $p > 0$, сумата им ще бъде безкрайност. Ако $p = 0$, сумата им ще е $0$. И в двата случая сумата не е $1$. 

\textbf{Пример 1 (Разпределение с "тежки опашки"):}
Можем да изберем:
\[ p_i = \frac{6}{\pi^2} \frac{1}{i^2}, \quad \text{за } i \ge 1 \]
Това е валидно вероятностно разпределение, защото от математическия анализ знаем, че $\sum_{i=1}^\infty \frac{1}{i^2} = \frac{\pi^2}{6}$, така че константата $\frac{6}{\pi^2}$ нормализира сумата до 1. Такива модели се ползват за екстремни събития, които макар и редки, имат значим ефект (тежки опашки).

\textbf{Пример 2 (Геометрично разпределение):}
Моделираме броя неуспехи преди първия успех. Нека $p \in (0, 1)$ и $q = 1 - p$. Тогава за $\Omega = \{0, 1, 2, 3, \dots\}$ полагаме:
\[ p_i = q^i p, \quad i \ge 0 \]
Това също сумира до 1 благодарение на формулата за сума на безкрайна геометрична прогресия:
\[ \sum_{i=0}^\infty q^i p = p \sum_{i=0}^\infty q^i = p \frac{1}{1-q} = \frac{p}{p} = 1 \]
Геометричното разпределение е интересно със свойството си за "безпаметност" – много подходящо за моделиране на живот на машинни елементи (но не и на хора, защото ние стареем).

\section{Геометрична вероятност}

Освен дискретното пространство, можем да имаме и непрекъснато. Нека $\Omega \subset \mathbb{R}^d$ е област в равнината (или пространството) с краен обем (или площ):
\[ \infty > \text{Vol}(\Omega) = |\Omega| = \int_{\Omega} dx \]
Тогава \textbf{геометричната вероятност} е дефинирана като равномерна вероятност върху това непрекъснато множество. За всяко подмножество $A \subseteq \Omega$ вероятността точката да попадне в $A$ е просто отношението на обемите:
\[ \mathbb{P}(A) := \frac{|A|}{|\Omega|} \]
Тук множеството от елементарни събития има неизброимо много точки. Ние не ползваме кардиналност (брой точки), а понятието обем (или площ). Всяка конкретна точка $x \in \Omega$ има обем $0$, следователно вероятността да се падне точно дадена конкретна точка е:
\[ \mathbb{P}(\{x\}) = \frac{0}{|\Omega|} = 0 \]
Това е фундаментално – при непрекъснатите разпределения вероятността за една конкретна точка е нула, въпреки че точката е възможен изход. Измерват се вероятности за попадане в области с положителен обем. Това свойство се използва например при методите Монте Карло, където генерираме случайни точки и оценяваме площи чрез съотношението на броя попаднали в $A$ точки към общия брой точки.

\section{Условна вероятност}

Накрая ще въведем едно много важно понятие, което отчита как се променят вероятностите при постъпване на нова информация. Да предположим, че имаме вероятностно пространство $(\Omega, \mathcal{A}, \mathbb{P})$. Наблюдаваме експеримента и научаваме, че се е случило събитието $A$ (където $\mathbb{P}(A) > 0$).

Тогава \textbf{условната вероятност} на събитие $B$ при дадено (или при условие) $A$ се дефинира като:
\[ \mathbb{P}(B \mid A) = \frac{\mathbb{P}(A \cap B)}{\mathbb{P}(A)} \]
Интуитивно, тъй като вече знаем, че сме попаднали в $A$, цялото ни пространство от възможни изходи се стеснява от $\Omega$ до $A$. Всички елементарни събития извън $A$ стават невъзможни и биват "изхвърлени". От събитието $B$ оцелява само тази част, която се пресича с $A$ (т.е. $A \cap B$).
За да получим коректна нова вероятност, която да сумира до 1, трябва да разделим (нормализираме) на вероятността на новото "пълно" пространство, което е $\mathbb{P}(A)$.

Например, в геометричната вероятност, ако първоначалната вероятност е била $\frac{|B|}{|\Omega|}$, след като знаем, че сме в $A$, новата вероятност става:
\[ \mathbb{P}(B \mid A) = \frac{|A \cap B| / |\Omega|}{|A| / |\Omega|} = \frac{|A \cap B|}{|A|} \]
Това преизчисляване е в основата на това как ние променяме своите очаквания въз основа на допълнителна информация. В следващата лекция ще разгледаме конкретни примери и приложения на условната вероятност.

\end{document}